\documentclass{beamer}

\usepackage{amsmath,amsfonts,amssymb,bm}
\usepackage{quantikz}
\usepackage{braket}
\usepackage{physics}
\usepackage{graphicx}

\usetheme{Madrid}

\title{Quantum Neural Networks}
\subtitle{Variational Circuits, Geometry, and Many-Body Connections}
\author{}
\date{}

\begin{document}

\frame{\titlepage}

%================================================
\section{Overview}
%================================================

\begin{frame}{What is a Quantum Neural Network?}
A QNN is a parameterized quantum circuit:
\[
\ket{\psi(x,\theta)} = U(\theta)U(x)\ket{0}
\]

Output:
\[
f(x,\theta) = \bra{\psi(x,\theta)} O \ket{\psi(x,\theta)}
\]

\begin{itemize}
\item $U(x)$: feature map
\item $U(\theta)$: trainable ansatz
\item $O$: observable
\end{itemize}
\end{frame}

%================================================
\section{Circuit Model}
%================================================

\begin{frame}{Basic QNN Circuit}
\[
\begin{quantikz}
\lstick{\ket{0}} & \gate{U(x)} & \gate{U(\theta)} & \meter{} \\
\lstick{\ket{0}} & \gate{U(x)} & \gate{U(\theta)} & \meter{}
\end{quantikz}
\]
\end{frame}

%------------------------------------------------

\begin{frame}{Example Ansatz}
\[
U(\theta) = \prod_i R_y(\theta_i)\prod_{i<j}\mathrm{CNOT}_{ij}
\]

\[
\begin{quantikz}
\lstick{\ket{0}} & \gate{R_y(\theta_1)} & \ctrl{1} & \qw \\
\lstick{\ket{0}} & \gate{R_y(\theta_2)} & \targ{} & \qw
\end{quantikz}
\]
\end{frame}

%================================================
\section{Pauli Expansion}
%================================================

\begin{frame}{Pauli Expansion of Ansatz}
\[
U(\theta) = e^{iH(\theta)}
\]

\[
H(\theta) = \sum_\alpha \theta_\alpha P_\alpha
\]

\begin{itemize}
\item $P_\alpha$: Pauli strings
\end{itemize}
\end{frame}

%------------------------------------------------

\begin{frame}{Example}
\[
H = \theta_1 Z_1 + \theta_2 Z_2 + \theta_3 Z_1Z_2
\]

\begin{itemize}
\item resembles spin Hamiltonian
\end{itemize}
\end{frame}

%================================================
\section{BCH Expansion}
%================================================

\begin{frame}{Effective Hamiltonian}
\[
U^\dagger O U = O + i[H,O] + \frac{i^2}{2}[H,[H,O]] + \cdots
\]
\end{frame}

%------------------------------------------------

\begin{frame}{Interpretation}
\begin{itemize}
\item nested commutators generate correlations
\item analogous to many-body expansions
\end{itemize}
\end{frame}

%================================================
\section{Expressivity and Entanglement}
%================================================

\begin{frame}{Entanglement}
\[
S(\rho_A) = -\mathrm{Tr}(\rho_A \log \rho_A)
\]
\end{frame}

%------------------------------------------------

\begin{frame}{Expressivity}
\begin{itemize}
\item low entanglement $\rightarrow$ simple functions
\item high entanglement $\rightarrow$ complex features
\end{itemize}
\end{frame}

%------------------------------------------------

\begin{frame}{Barren Plateaus}
\[
\mathrm{Var}(\nabla C) \sim e^{-n}
\]
\end{frame}

%================================================
\section{Quantum Fisher Information}
%================================================

\begin{frame}{QFI}
\[
F_{ij} = 4\,\mathrm{Re}(\langle \partial_i \psi | \partial_j \psi\rangle)
\]
\end{frame}

%------------------------------------------------

\begin{frame}{Natural Gradient}
\[
\dot{\theta} = F^{-1} \nabla C
\]
\end{frame}

%================================================
\section{Time-Dependent Variational Principle}
%================================================

\begin{frame}{TDVP}
\[
\delta \| (i\partial_t - H)|\psi\rangle \| = 0
\]
\end{frame}

%------------------------------------------------

\begin{frame}{Equation of Motion}
\[
\sum_j F_{ij} \dot{\theta}_j = C_i
\]

\[
C_i = \mathrm{Im}\langle \partial_i \psi | H | \psi\rangle
\]
\end{frame}

%================================================
\section{Connection to ADAPT-VQE}
%================================================

\begin{frame}{ADAPT-VQE}
\[
\langle [H,A_k]\rangle
\]

\begin{itemize}
\item selects operators dynamically
\end{itemize}
\end{frame}

%------------------------------------------------

\begin{frame}{Relation to QNN}
\begin{itemize}
\item QNN: fixed ansatz
\item ADAPT: growing ansatz
\end{itemize}
\end{frame}

%================================================
\section{Training}
%================================================

\begin{frame}{Parameter Shift}
\[
\frac{\partial f}{\partial \theta}
=
\frac{1}{2}(f(\theta+\frac{\pi}{2})-f(\theta-\frac{\pi}{2}))
\]
\end{frame}

%================================================
\section{Example Circuit}
%================================================

\begin{frame}{QNN Example}
\[
\begin{quantikz}
\lstick{\ket{0}} & \gate{R_z(x_1)} & \gate{R_y(\theta_1)} & \ctrl{1} & \qw \\
\lstick{\ket{0}} & \gate{R_z(x_2)} & \gate{R_y(\theta_2)} & \targ{} & \qw
\end{quantikz}
\]
\end{frame}

%================================================
\section{Physics Interpretation}
%================================================

\begin{frame}{QNN as Quantum Dynamics}
\[
|\psi\rangle = e^{-iH(\theta)}|\psi(x)\rangle
\]

\begin{itemize}
\item learning = tuning Hamiltonian
\end{itemize}
\end{frame}

%------------------------------------------------

\begin{frame}{Connection to Many-Body Theory}
\begin{itemize}
\item BCH $\leftrightarrow$ perturbation theory
\item QFI $\leftrightarrow$ metric tensor
\item TDVP $\leftrightarrow$ classical dynamics
\end{itemize}
\end{frame}

%================================================
\section{Summary}
%================================================

\begin{frame}{Summary}
\begin{itemize}
\item QNN = variational quantum circuit
\item learning = optimization in Hilbert space
\item geometry governed by QFI
\item strong connection to many-body physics
\end{itemize}
\end{frame}

\end{document}
